% Generated draft: requires mathematical and source review.
\documentclass[11pt]{article}
\usepackage[margin=25mm]{geometry}
\usepackage{fontspec}
\setmainfont{Latin Modern Roman}
\usepackage{amsmath,amssymb,mathtools,unicode-math}
\setmathfont{Latin Modern Math}
\usepackage{enumitem,needspace,xurl,hyperref}
\setlength{\parindent}{0pt}
\setlength{\parskip}{5pt}
\setlength{\emergencystretch}{3em}
\newcommand{\N}{\mathbb{N}}
\newcommand{\Z}{\mathbb{Z}}
\newcommand{\Q}{\mathbb{Q}}
\newcommand{\R}{\mathbb{R}}
\newcommand{\C}{\mathbb{C}}
\newcommand{\F}{\mathbb{F}}
\newcommand{\A}{\mathbb{A}}
\newcommand{\PP}{\mathbb{P}}
\newcommand{\E}{\mathbb{E}}
\newcommand{\eps}{\varepsilon}
\newcommand{\dd}{\,\mathrm{d}}
\begin{document}
\section*{UnsolvedMath ID: 30004709}
\textbf{Source problem code:} catalogue entry
\subsection{Definitions and mathematical statement}
Use the standard normalization that a hyperbolic metric has sectional curvature \(-1\). For a fixed integer \(n\geq 2\) and a real number \(v>0\), let
\[
\mathcal H_n(v)=\{M: M\text{ is a connected, oriented, closed hyperbolic }n\text{-manifold},\ \operatorname{Vol}(M)\leq v\}/\text{isometry}.
\]
The range \(n\geq2\) is the conventional interpretation of the dimension parameter; the supplied record does not explicitly specify its range. Here closed means compact without boundary.

For such a manifold \(M\), define its systole by
\[
\operatorname{sys}(M)=\min\{\ell(\gamma):\gamma\text{ is a nonconstant closed geodesic in }M\}.
\]
Under the convention that geodesics are counted without orientation or parametrization, define the kissing number by
\[
\operatorname{kiss}(M)=\#\{\gamma:\gamma\text{ is a geometrically distinct closed geodesic of length }\operatorname{sys}(M)\}.
\]
The supplied record does not specify the counting convention for the kissing number, so statements involving this definition are conditional on that convention.

For the nonnegative Laplace--Beltrami operator \(\Delta_M\) on scalar functions, let \(\lambda_1(M)\) be its first strictly positive eigenvalue, and put
\[
m_1(M)=\dim_{\R}\{f\in C^\infty(M,\R):\Delta_M f=\lambda_1(M)f\}.
\]
For each invariant
\[
I\in\{\operatorname{sys},\operatorname{kiss},\lambda_1,m_1\},
\]
and each \(v\) for which \(\mathcal H_n(v)\neq\varnothing\), set
\[
F_{n,I}(v)=\sup_{M\in\mathcal H_n(v)} I(M).
\]
Using a supremum does not assert finiteness or attainment.

The problem is to determine these extremal functions and the manifolds attaining their values for low volume bounds, and to determine their asymptotic behavior as \(v\to\infty\), with \(n\) fixed. The source does not specify a numerical meaning of low volume or the desired strength of the asymptotic description. This formulation treats the four invariants as separate optimization problems; the supplied record does not explicitly resolve whether any simultaneous optimization is intended.

In particular, let \(\mathcal G_2\) be the class of connected, oriented, closed hyperbolic surfaces of genus two, and let \(B\) denote the Bolza surface with its hyperbolic metric. Are both assertions
\[
\lambda_1(S)\leq\lambda_1(B)\quad\text{for every }S\in\mathcal G_2,
\qquad
m_1(S)\leq m_1(B)\quad\text{for every }S\in\mathcal G_2
\]
true? These are separate extremality assertions and do not assert uniqueness. Under the curvature normalization above, every surface in \(\mathcal G_2\) has area \(4\pi\).

The supplied dated assessment labels the problem open, describes Bolza optimality for the systolic quantities in the stated context as known, and describes its expected optimality for the first eigenvalue and multiplicity, together with the requested asymptotics, as unproved. These are supplied curation claims, not independently verified mathematical conclusions.
\subsection{Short English statement}
For each fixed dimension, maximize the systole, number of shortest closed geodesics, first positive Laplace eigenvalue, and its multiplicity over closed oriented hyperbolic manifolds with volume at most \(v\). Determine low-volume extremizers and large-\(v\) behavior. In genus two, ask whether the Bolza surface separately maximizes the first positive eigenvalue and its multiplicity. Counting conventions and the intended precision of the asymptotics require review.
\subsection{Sources}
\begin{itemize}
\item Local UnsolvedMath registry: permanent ID \(30004709\), exact source problem code \texttt{catalogue entry}, title ``Low-Volume and Asymptotic Hyperbolic-Invariant Maximizers,'' in the Oberwolfach Reports Open Problems collection. No individual proposer is supplied; the recorded proposed year is 2021.
\item Supplied report citation: \emph{Enumerative Geometry of Surfaces (2022)}, Oberwolfach Reports, pp.\ 1463--1518, DOI \texttt{10.4171/owr/2021/28}. \url{https://doi.org/10.4171/owr/2021/28}. The supplied citation uses 2022, while the DOI and recorded proposed year refer to 2021; these details are retained without reconciliation.
\item Supplied literature source, labelled ``Fortier Bourque--Petri'': \url{https://arxiv.org/abs/2302.02540}. No paper title or specific supporting result is supplied.
\item Literature assessment checked at \texttt{2026-08-22}: supplied dated curation, not independently verified. It labels the status open, states that the Bolza surface is known to optimize the systolic quantities in the stated genus-two context, and describes the eigenvalue and multiplicity optimality questions and requested asymptotics as unproved. The background explicitly requests independent verification of time-sensitive status and cited claims.
\item The registry labels the statement \texttt{corrected\_verified} and reports that its standalone formulation and hypotheses were checked against the official Oberwolfach report. This is supplied verification metadata; no source was independently visited or checked in preparing this draft.
\end{itemize}
\end{document}
